\section{Digital signatures}\label{sec:prelims-signatures}

Digital signature schemes enable signers to produce publicly verifiable
signatures on messages.

\paragraph{Syntax.} A digital signature scheme consists of the following
algorithms:

\begin{description}

  \item[$(\pk,\sk) \sample \DSgen(\secparam)$ : ] The key generation algorithm
    outputs verification key $\pk$ and signing key $\sk$.

  \item[$\signature \sample \DSsign(\sk, \msg)$ : ] The signing algorithm signs
    message $\msg$ using secret key $\sk$, giving output $\signature$.

  \item[$\{0,1\} \leftarrow \DSverify(\pk,\msg,\signature)$ : ] The
    verification algorithm verifies signature $\signature$ on message $\msg$
    using public key $\pk$, giving output $0$ or $1$.

\end{description}

Digital signature schemes satisfy correctness and unforgeability, one form of
which we define below. Again, we use standard syntax and security notions for
digital signatures~\cite{KatLin14}.

\paragraph{Correctness.} A digital signature scheme is correct if, for every
key pair $(\pk,\sk)$ generated by $\DSgen(\secparam)$ and every message $\msg$,
the following holds:
%
\begin{equation*}
%
  \DSverify(\pk, \msg, \DSsign(\sk, \msg)) = 1.
%
\end{equation*}
%

\paragraph{Unforgeability.} We define existential unforgeability under chosen
message attack (EUF-CMA) using the following security experiment.

\begin{experimentbox}{Existential unforgeability under chosen message
  attack}{euf-cma}

  \begin{enumerate}

    \item The challenger generates $(\pk,\sk) \sample \DSgen(\secparam)$ and
      sends $\pk$ to $\adv$.

    \item The adversary is given oracle access to $\DSsigno$, which on input a
      message $\msg$ returns $\signature = \DSsign(\sk, \msg)$ and stores
      $\msg$ in the set $\mssg$.

    \item The adversary outputs $(\msg^\star, \signature^\star)$.

  \end{enumerate}

  $\adv$ wins the experiment if the following hold:
  %
  \begin{equation*}
  %
    \DSverify(\pk,\msg^\star,\signature^\star) = 1 \ \wedge\
    \msg^\star \notin \mssg.
  %
  \end{equation*}

\end{experimentbox}

The advantage of $\adv$ is defined as follows:
%
\begin{equation*}
%
  \Adv^{\mathsf{euf\text{-}cma}}_{\DS, \ \adv}(\secparam) = \Pr[ \ \adv\text{
  wins} \ ].
%
\end{equation*}
%
We say the digital signature scheme $\DS$ is EUF-CMA secure if this advantage
is small for every admissible adversary $\adv$.